\chapter{卷积网络：共享局部计算}
\label{ch:cnn}

卷积神经网络（Convolutional Neural Network, CNN）把局部连接与参数共享结合起来，
适合处理具有网格结构的输入。理解CNN时，先跟踪形状、一次局部计算和参数数量，
再理解多层组合怎样形成表示。

\begin{keyidea}[同一个核，在多个位置重复使用]
一个卷积核在不同空间位置使用同一组权重。
这能减少参数，并表达对局部模式重复出现的偏好；它不保证对所有位置变化完全不敏感。
\end{keyidea}

\section{张量形状：批量、通道与空间}
\label{sec:cnn-tensor}
深度学习程序中的张量（Tensor）通常指多维数组。
图像批量常用NCHW或NHWC布局；字母顺序决定轴的含义，不能只看数组有几维。
本章统一用$B,C,H,W$表示批量、通道、高、宽。

\begin{figure}[htbp]
\centering\begin{tikzpicture}[node distance=7mm]
\node[box,text width=24mm] (b) {$B=8$\\批量样本};
\node[box,text width=24mm,right=of b] (c) {$C=3$\\颜色通道};
\node[box,text width=24mm,right=of c] (h) {$H=32$\\高度};
\node[box,text width=24mm,right=of h] (w) {$W=32$\\宽度};
\node[lab,below=10mm of $(c)!0.5!(h)$] {NCHW布局示例：$X\in\R^{8\times3\times32\times32}$};
\end{tikzpicture}

\caption{8张32$\times$32的RGB图像，按通道在前布局存储。每个通道不是一条独立样本。}
\par\smallskip
\readfigure{共有$8\times3\times32\times32=24576$个数值。
卷积通常在每张图像上独立进行，跨输入通道加权求和，不把批量轴当空间轴滑动。}
\end{figure}

标量、向量、矩阵分别可以作为0阶、1阶、2阶数组；这里的“阶/轴数”不要与矩阵秩或特征数混用。
改变轴顺序需要相应置换，单纯reshape通常不能代替布局转换。

\quick{每层都记录输入、输出的轴顺序与形状，比仅画一个三维方块更容易检查实现。}

\section{一次卷积怎样得到一个数？}
\label{sec:cnn-convolution}
常见深度学习库称为“卷积”的操作实际按不翻转核的互相关（Cross-correlation）实现。
单通道、步幅1、无填充时，$Y_{ij}=b+\sum_{u,v}K_{uv}X_{i+u,j+v}$。
数学上的离散卷积常还包含核翻转，读公式时需声明约定。

\begin{figure}[htbp]
\centering\begin{tikzpicture}[x=8mm,y=8mm]
\foreach \r/\a/\b/\c/\d in {0/1/2/0/1,1/0/1/3/2,2/2/0/1/1,3/1/2/1/0}{
 \foreach \j/\v in {0/\a,1/\b,2/\c,3/\d}{
  \ifnum\r<2\ifnum\j<2\def\cc{Pale}\else\def\cc{white}\fi\else\def\cc{white}\fi
  \draw[Teal!40,fill=\cc] (\j,-\r) rectangle ++(1,-1);
  \node at (\j+.5,-\r-.5) {$\v$};}}
\node[lab] at (2,1) {输入$X$};
\foreach \r/\a/\b in {0/1/0,1/0/-1}{\foreach \j/\v in {0/\a,1/\b}{
 \draw[Teal,fill=Pale] (6+\j,-\r-1) rectangle ++(1,-1);
 \node at (6.5+\j,-\r-1.5) {$\v$};}}
\node[lab] at (7,1) {共享核$K$};
\draw[flow] (8.5,-2)--(10,-2);
\foreach \r/\a/\b/\c in {0/0/-1/-2,1/0/0/2,2/0/-1/1}{\foreach \j/\v in {0/\a,1/\b,2/\c}{
 \draw[Teal!40,fill=Pale] (11+\j,-\r-.5) rectangle ++(1,-1);
 \node at (11.5+\j,-\r-1) {$\v$};}}
\node[lab] at (12.5,1) {输出$Y$};
\draw[Amber,thick] (11,-.5) rectangle ++(1,-1);
\node[lab] at (7,-5) {左上角：$1\times1+2\times0+0\times0+1\times(-1)=0$};
\end{tikzpicture}

\caption{4$\times$4输入、2$\times$2核，偏置0；每次移动一格，得到3$\times$3输出。}
\par\smallskip
\readfigure{核左上位置乘1、右下位置乘$-1$，其余乘0。
在输入左上2$\times$2块计算得0；向右一格得到$2-3=-1$。
输出每个位置都复用同一组四个权重。}
\end{figure}

核权重通常由损失与反向传播学习，不能一概说每个核固定是边缘检测器。
叠加非线性使多层网络能表达更复杂关系；若只叠加线性或仿射层而没有非线性，
整体仍是线性或仿射映射。

\begin{pitfall}[特征图不是原图的复制]
输出位置表示某个局部加权计算的结果。空间大小、通道数与原输入可能不同，
数值也不必在原像素范围内。
\end{pitfall}

\section{通道与参数共享：参数数怎样算？}
\label{sec:cnn-parameters}
普通非分组二维卷积中，每个输出通道都有覆盖所有输入通道的核：
\[
 Y_{o,i,j}=b_o+\sum_{c=1}^{C_{\mathrm{in}}}\sum_{u,v}K_{o,c,u,v}X_{c,i+u,j+v}.
\]
$k_h\times k_w$核、每输出通道一个偏置时，参数数为
\[
 C_{\mathrm{out}}(C_{\mathrm{in}}k_hk_w+1).
\]
若不使用偏置，去掉最后的$C_{\mathrm{out}}$；分组卷积要按分组方式另算。

\begin{center}
\begin{tabularx}{\textwidth}{P{52mm}P{28mm}Y}
\toprule
相同3$\times$3输出的教学比较 & 参数数 & 连接方式\\
\midrule
4$\times$4单通道输入，2$\times$2卷积 & $4+1=5$ & 9个位置共享一个核与偏置。\\
同样局部窗口，但各位置不共享 & $9(4+1)=45$ & 每位置单独学习。\\
完全连接到9个输出单元 & $9(16+1)=153$ & 每单元连接全部16输入。\\
\bottomrule
\end{tabularx}
\end{center}

较大的例子：3输入通道、16输出通道、3$\times$3核，共$16(27+1)=448$个参数。
同一层参数数不随批量或图像宽高增加而增加，但计算量和激活内存会增加。

\quick{通道汇总、空间滑动和批量独立，是三种不同的轴操作。}

\section{步幅、填充与输出大小}
\label{sec:cnn-shape}
对某一空间轴，输入长度$L$、两侧各填充$p$、核大小$k$、步幅$s$、膨胀率$d$，
常用输出长度为
\[
 L_{\mathrm{out}}=\left\lfloor\frac{L+2p-d(k-1)-1}{s}\right\rfloor+1.
\]
有效核跨度为$d(k-1)+1$。公式假定参数有效、使用通常的向下取整规则，
某些自动padding或池化的ceil模式需要另按接口约定。

\begin{center}
\begin{tabular}{lrrrrrr}
\toprule
例子 & $L$ & $k$ & $p$ & $s$ & $d$ & 输出\\
\midrule
无填充小例子 & 4 & 2 & 0 & 1 & 1 & 3\\
保持32长度 & 32 & 3 & 1 & 1 & 1 & 32\\
步幅2降采样 & 7 & 3 & 1 & 2 & 1 & 4\\
膨胀核示例 & 7 & 3 & 0 & 1 & 2 & 3\\
\bottomrule
\end{tabular}
\end{center}

Padding在边缘提供约定的额外数值，例如0；它不是自动补回真实图像信息。
Stride跳过部分位置，Dilation拉开核采样点间隔；二者对输出大小与感受野的作用不同。

\begin{pitfall}[“same”不总等于宽高完全不变]
一些接口在步幅1下维持大小，在步幅大于1时按特定规则产生约$L/s$的输出。
使用文字模式时，应查接口的填充与非对称边界约定。
\end{pitfall}

\section{池化与理论感受野}
\label{sec:cnn-pooling}
池化（Pooling）按局部窗口汇总数值，常见最大池化和平均池化通常没有可学习权重。
它们可减少空间大小，但也可能丢失位置或细节信息。

\begin{figure}[htbp]
\centering\begin{tikzpicture}[x=8mm,y=8mm]
\foreach \r/\a/\b/\c/\d in {0/1/3/2/0,1/4/6/1/5,2/2/1/8/7,3/0/3/4/9}{\foreach \j/\v in {0/\a,1/\b,2/\c,3/\d}{
 \draw[Teal!40,fill=Pale] (\j,-\r) rectangle ++(1,-1);\node at (\j+.5,-\r-.5) {$\v$};}}
\draw[Amber,thick] (0,0) rectangle (2,-2);
\node[lab] at (2,1) {输入4$\times$4};
\draw[flow] (4.5,-2)--(6,-2);
\foreach \r/\a/\b in {0/6/5,1/3/9}{\foreach \j/\v in {0/\a,1/\b}{
 \draw[Teal!40,fill=Pale] (7+\j,-\r-1) rectangle ++(1,-1);\node at (7.5+\j,-\r-1.5) {$\v$};}}
\node[lab] at (8,1) {最大池化};
\foreach \r/\a/\b in {0/3.5/2,1/1.5/7}{\foreach \j/\v in {0/\a,1/\b}{
 \draw[Teal!40,fill=Pale] (12+\j,-\r-1) rectangle ++(1,-1);\node at (12.5+\j,-\r-1.5) {$\v$};}}
\node[lab] at (13,1) {平均池化};
\node[lab] at (7,-5) {窗口2$\times$2，步幅2，无填充；两种池化是不同选择};
\end{tikzpicture}

\caption{两个独立的2$\times$2池化例子；最大与平均不是要依次串联的两步。}
\par\smallskip
\readfigure{左上块$[1,3;4,6]$最大值为6，均值为3.5。
最大池化的反向梯度流向取得最大值的位置；出现相等最大值时需按实现约定处理。}
\end{figure}

一个输出所能依赖的输入区域称感受野（Receptive Field）。
令初始跨度$r_0=1$、相邻位置间距$j_0=1$，则
\[
 r_l=r_{l-1}+[d_l(k_l-1)]j_{l-1},\qquad j_l=j_{l-1}s_l.
\]
3核步幅1卷积后，$r=3,j=1$；2核步幅2池化后，$r=4,j=2$；
再接3核步幅1卷积，$r=8,j=2$。这是理论依赖跨度，不表示该区域内每个像素影响完全相同。

\quick{共享卷积常带来平移等变倾向；步幅、填充和边界会影响它，池化也不保证任意平移不变。}

\section{残差连接：为变换加一条直接路径}
\label{sec:cnn-residual}
残差块（Residual Block）把一个变换$F(x)$与输入相加：$y=x+F(x)$。
当$F(x)=0$时，可以直接实现恒等映射；这为深层网络的信息与梯度提供另一条路径。

\begin{figure}[htbp]
\centering\begin{tikzpicture}[node distance=12mm]
\node[box,text width=21mm] (x) {输入$x$};
\node[box,text width=36mm,right=of x] (f) {若干变换\\$F(x)$};
\node[circle,draw=Teal,fill=Pale,minimum size=9mm,right=of f] (add) {$+$};
\node[box,text width=30mm,right=of add] (y) {$y=x+F(x)$};
\draw[flow] (x)--(f);\draw[flow] (f)--(add);\draw[flow] (add)--(y);
\draw[flow] (x.north)--++(0,12mm)-|node[pos=.3,above,lab]{恒等支路}(add.north);
\end{tikzpicture}

\caption{最简残差结构示意，未规定卷积、归一化和激活的具体顺序。}
\par\smallskip
\readfigure{两条支路的输出形状必须能逐元素相加。
若通道数或空间大小变化，可用相应投影/下采样支路$P(x)$，成为$y=P(x)+F(x)$。}
\end{figure}

在形状一致、可微设置中，局部Jacobian为$I+\partial F/\partial x$。
这说明直接路径的存在，但不意味着梯度永不消失或训练必定成功。
不同ResNet变体的归一化与激活位置不同，应按具体结构计算。

\begin{pitfall}[相加与拼接不同]
相加要求对应位置具有兼容形状，输出通道数通常不增加。
拼接（Concatenation）沿指定轴保留两组特征，会改变该轴长度。
不要因为图中都是两条箭头汇合，就把两种操作视作同一个运算。
\end{pitfall}

\section{归一化层：统计量在哪些轴上计算？}
\label{sec:cnn-normalization}
Batch Normalization（BN）训练时按约定轴计算均值与方差，再进行可学习的缩放与平移：
\[
 \hat x=\frac{x-\mu}{\sqrt{v+\epsilon}},\qquad y=\gamma\hat x+\beta.
\]
CNN中常为每个通道在批量和空间位置上汇总。$\epsilon>0$用于数值稳定，
$\gamma,\beta$为可学习参数，统计量不是同一种参数。

\begin{center}
\begin{tabularx}{\textwidth}{P{33mm}YY}
\toprule
层/阶段 & 常见统计范围 & 预测时的考虑\\
\midrule
BN训练 & 同通道的批量与空间位置。 & 训练批量影响统计，常维护滑动估计。\\
BN推理 & 通常使用保存的滑动统计。 & 需切换到适当模式，不能默认仍用当前批量。\\
Layer Normalization & 单条记录约定的特征轴。 & 通常训练/推理都按该条输入计算，无批量滑动统计。\\
\bottomrule
\end{tabularx}
\end{center}

简单两值$x=[1,3]$，$\mu=2,v=1$，$\gamma=1,\beta=0$时，
输出为$[-1,1]/\sqrt{1+\epsilon}$。
方差估计和滑动更新的细节取决于框架，不能把所有实现的分母约定混为一谈。

\quick{训练/推理模式与是否记录梯度是不同开关；归一化也不会自动把分布变成正态。}

\section{数据增强：保持任务语义的变化}
\label{sec:cnn-augmentation}
数据增强（Data Augmentation）在训练中构造合理变换，让模型学习任务所允许的变化。
它是建模假设，不是“加入任意噪声都能大幅改善表现”。

\begin{center}
\begin{tabularx}{\textwidth}{P{33mm}Y}
\toprule
操作 & 必须同时考虑\\
\midrule
裁剪、翻转、旋转 & 是否改变标签？方向识别、数字或医学左右侧任务可能不适合某些变换。\\
亮度或颜色扰动 & 是否仍处于真实采集的合理范围？颜色本身可能是目标信号。\\
检测与分割变换 & 框、掩码与图像使用同一几何变换。\\
训练数据派生记录 & 先按原对象/时间划分，再在训练内增强，避免同源副本跨验证测试。\\
\bottomrule
\end{tabularx}
\end{center}

增强样本通常与原记录相关，不能当作独立新增的真实对象来估计有效样本量。
验证与测试使用约定的确定性处理；测试时增强若使用，应作为完整评估方案明确报告。

\begin{pitfall}[共享权重不保证识别所有位置]
参数共享使相同局部计算在多处可用；识别完整对象仍依赖训练分布、上下文与后续层。
CNN对旋转、尺度和任意位置变化的鲁棒性不是自动得到的保证。
\end{pitfall}

\section{一个完整小网络：形状与参数对照}
\label{sec:cnn-example}
下面是用于说明计算的小型10类图像分类器，不代表对任何任务的最佳架构。
卷积采用3核、步幅1、填充1，池化采用2核、步幅2。

\begin{figure}[htbp]
\centering\begin{tikzpicture}[node distance=10mm and 13mm]
\node[box,text width=33mm] (input) {输入\\$8\times3\times32\times32$};
\node[box,text width=33mm,right=of input] (conv) {Conv3$\times$3+ReLU\\$8\times16\times32\times32$};
\node[box,text width=33mm,right=of conv] (pool) {MaxPool2$\times$2\\$8\times16\times16\times16$};
\node[box,text width=33mm,below=of pool] (gap) {全局平均池化\\$8\times16$};
\node[box,text width=33mm,left=of gap] (head) {线性分类头\\$8\times10$};
\node[box,text width=33mm,left=of head] (loss) {10类logits\\训练用交叉熵};
\draw[flow] (input)--(conv);\draw[flow] (conv)--(pool);\draw[flow] (pool)--(gap);
\draw[flow] (gap)--(head);\draw[flow] (head)--(loss);
\end{tikzpicture}

\caption{蛇形箭头按顺序连接，批量轴8始终保留；最后输出尚未归一化的logits。}
\par\smallskip
\readfigure{卷积有448个参数；线性头16输入、10输出，含偏置共170个参数。
ReLU、池化和全局平均池化无可学习参数，本例总计618个参数。}
\end{figure}

全局平均池化（Global Average Pooling）对每通道的空间位置取平均，
把$8\times16\times16\times16$变为$8\times16$，不对批量或通道轴求平均。
训练时常把logits直接交给包含稳定对数运算的交叉熵接口，避免重复Softmax。

\subsection{自测与回看}
图像宽高翻倍，卷积参数会自动翻四倍吗？\textbf{共享核参数不变，计算与激活内存会增大。}
张量4轴是否代表四维特征空间的四个特征？\textbf{不是，必须看每轴含义。}
残差相加是否与通道拼接相同？\textbf{不同。}

\begingroup\small
\noindent\textbf{参考：}Goodfellow等，\textit{Deep Learning}（2016）第9章；
He等，\textit{Deep Residual Learning for Image Recognition}（2016）；PyTorch Conv2d、Pooling与Normalization文档。
图和例子为独立构造。
\par\endgroup
